\section{Discussion and Conclusion}
\begin{frame}[t]{Conclusions and Organizational Implications}
\small
\begin{block}{AI changes the resources available to a worker}
AI-assisted individuals can achieve output quality associated with cross-functional teams. AI also broadens the orientation of their proposals and improves the emotional experience of the task.
\end{block}
\begin{block}{Human judgment remains valuable}
The process results distinguish the ability to generate strong ideas from the ability to select the best one. Work design can combine AI-supported exploration with deliberate evaluation.
\end{block}
\begin{alertblock}{Questions for organizational design}
How should tasks be allocated between individuals and teams? Which review practices improve selection? How does repeated AI use affect the development of expertise?
\end{alertblock}

\end{frame}

\begin{frame}[t]{Reference}
\small

\begin{block}{Article}
Dell'Acqua, F., Ayoubi, C., Lifshitz, H., Sadun, R., Mollick, E., Mollick, L., Han, Y., Goldman, J., Nair, H., Taub, S., and Lakhani, K. R. (2026).
\medskip\par
\emph{The Cybernetic Teammate: A Field Experiment on Generative AI and Teamwork.}
\medskip\par
Organization Science, 37(4), 1217--1242.
\medskip\par
\href{https://doi.org/10.1287/orsc.2025.20702}{doi.org/10.1287/orsc.2025.20702}
\end{block}

\end{frame}

